% leanmap — mathematical foundations
% Compile: pdflatex leanmap && bibtex leanmap && pdflatex leanmap && pdflatex leanmap
% Or (no bib needed for first pass): pdflatex leanmap
\documentclass[11pt,a4paper]{article}

\usepackage[margin=1in]{geometry}
\usepackage{amsmath,amssymb,amsthm,mathtools}
\usepackage{bm}
\usepackage{booktabs}
\usepackage{enumitem}
\usepackage{hyperref}
\usepackage{microtype}
\usepackage{xcolor}
\usepackage[numbers,sort&compress]{natbib}

\hypersetup{
  colorlinks=true,
  linkcolor=blue!50!black,
  citecolor=blue!50!black,
  urlcolor=blue!40!black,
}

\newtheorem{definition}{Definition}
\newtheorem{proposition}{Proposition}
\newtheorem{remark}{Remark}
\newtheorem{assumption}{Assumption}

\newcommand{\R}{\mathbb{R}}
\newcommand{\E}{\mathbb{E}}
\newcommand{\norm}[1]{\left\|#1\right\|}
\newcommand{\abs}[1]{\left|#1\right|}
\DeclareMathOperator*{\argmin}{arg\,min}
\newcommand{\inner}[2]{\left\langle #1,#2\right\rangle}
\newcommand{\clamp}{\mathrm{clamp}}
\newcommand{\softplus}{\mathrm{softplus}}
\newcommand{\softmax}{\mathrm{softmax}}
\newcommand{\GELU}{\mathrm{GELU}}
\newcommand{\LN}{\mathrm{LN}}
\newcommand{\ind}{\mathbf{1}}
\newcommand{\loss}{\mathcal{L}}

\title{Mathematical Foundations of leanmap\\
\large Parametric Landmark-Conditioned Neighbour Embedding\\
with Multi-Scale Graph Pyramids and Metric Regularisers}
\author{leanmap technical note}
\date{\today}

\begin{document}
\maketitle

\begin{abstract}
leanmap (\emph{LEarn ANother MAPping}) is a parametric neighbour-embedding
method: a feed-forward map \(f_\theta:\R^D\to\R^d\) is trained once against a
fuzzy simplicial graph and then evaluates new points by a single network
forward pass. The neighbour graph is discarded after training; the saved
artefact holds network weights, landmarks, and conformal calibration scores.
This note records the mathematics implemented in the reference package
(\texttt{PLANE} / \texttt{PLANEConfig}): graph construction (including an
\(\varepsilon\)-net and a Galerkin pyramid), the FiLM-conditioned encoder,
the composite training objective (fuzzy cross-entropy, ordinal triplets,
landmark load-balancing, Gram-domain local rigidity, and geodesic Procrustes /
stress), the \((a,b)\) curve fit that
implements \texttt{min\_dist}, ramp schedules, evaluation statistics, and
conformal out-of-distribution scores.

Section~\ref{sec:measurements} reports the seeded evidence behind the shipping
defaults. Every number there is a mean over three seeds with the spread given,
and a difference is called resolved only when it exceeds twice the pooled seed
spread. Several claims made in earlier versions of this note do not survive
that test and are withdrawn where they appear.
\end{abstract}

\tableofcontents
\newpage

%==============================================================================
\section{Setup and notation}
%==============================================================================

\subsection{Problem}
\label{sec:problem}

Given a finite sample
\[
X = \{x_1,\ldots,x_N\}\subset\R^D
\]
and a distance function \(d_X:\R^D\times\R^D\to[0,\infty)\) (the
\emph{ambient metric}), leanmap constructs a parametric embedding
\[
f_\theta:\R^D\to\R^d,\qquad d=2\text{ by default},
\]
such that local neighbourhood structure in \(X\) is preserved in the image
\(Z=\{z_i=f_\theta(x_i)\}\), global geodesic structure is approximately
retained when requested, and \(f_\theta\) remains well-defined for previously
unseen points.

Unlike classical UMAP~\citep{mcinnes2018umap}, which stores embedding
coordinates as free variables, leanmap stores a network. After training the
fuzzy graph is discarded; inference is a forward pass plus a cover score
(Section~\ref{sec:conformal}), costing \(3.4\,\mu\)s per point against
\(781\,\mu\)s for \texttt{UMAP.transform} on the same data
(Section~\ref{sec:inference}). What that buys, and what it gives up, is stated
precisely in Remark~\ref{rem:notindexfree}: the neighbour search is replaced by
a fixed-size computation rather than eliminated.

\subsection{Symbols}
\label{sec:symbols}

\begin{center}
\small
\begin{tabular}{ll}
\toprule
Symbol & Meaning \\
\midrule
\(N,D\) & number of training points; ambient dimension \\
\(R\) & number of \(\varepsilon\)-net representatives (\(R\le N\)) \\
\(L\) & number of landmarks / anchors \\
\(d\) & embedding dimension (\texttt{d\_out}) \\
\(k\) & neighbourhood size (\texttt{n\_neighbors}) \\
\(x_i,z_i\) & ambient point; embedding \(z_i=f_\theta(x_i)\) \\
\(M_\ell\in\R^D\) & landmark coordinates, \(\ell=1,\ldots,L\) \\
\(D_{i\ell},a_{i\ell}\) & point--landmark distance and affinity \\
\(\tau_\ell\) & per-landmark temperature \\
\(\rho_i,\sigma_i\) & local connectivity offset; bandwidth \\
\(p_{ij},P^{\mathrm{sym}}\) & directed memberships; symmetrised fuzzy graph \\
\(w_{ij}\) & edge weight in a training batch \\
\(a,b\) & fitted low-dimensional kernel parameters \\
\(\gamma_k,\beta_k\) & FiLM scale / shift at layer \(k\) \\
\(t\in[0,1]\) & training fraction (for ramps) \\
\(g_{\ell\ell'}\) & landmark--landmark graph geodesic \\
\bottomrule
\end{tabular}
\end{center}

Distances used by the graph and losses are measured in units of a fitted
\emph{natural scale}: the median \(k\)-NN distance on a subsample of \(X\),
so that neighbourhood radii are \(\mathcal{O}(1)\) regardless of the raw
units of \(d_X\).

%==============================================================================
\section{Overview of the method}
\label{sec:overview}
%==============================================================================

Training proceeds in a fixed order:

\begin{enumerate}[leftmargin=*]
\item Hold out a conformal calibration split from raw \(X\)
  (fraction \(\texttt{CALIB\_FRAC}=0.05\), capped by \texttt{calib\_max}).
\item Fit encoder normalisation \((\mu,\sigma)\) and the metric natural scale.
\item Select landmarks \(M\) and build a multi-scale fuzzy graph pyramid
  (Section~\ref{sec:graph}).
\item Instantiate the FiLM encoder \(f_\theta\) (Section~\ref{sec:model}).
\item Fit kernel parameters \((a,b)\) from \texttt{min\_dist}
  (Section~\ref{sec:ab}).
\item Optimise the composite loss (Section~\ref{sec:loss}) by AdamW.
\item Calibrate conformal cover scores on the held-out split
  (Section~\ref{sec:conformal}).
\end{enumerate}

The per-step training objective is
\begin{equation}
\label{eq:total}
\begin{aligned}
\loss
&=
\loss_{\mathrm{geom}}
+ \lambda_{\mathrm{ord}}\,\loss_{\mathrm{ord}}
+ \lambda_{\mathrm{lm}}\,\loss_{\mathrm{lm}} \\
&\quad
+ \lambda_{\mathrm{frame}}\,r_{\mathrm{frame}}(t)\,\loss_{\mathrm{frame}}
+ \lambda_{\mathrm{geo}}\,r_{\mathrm{geo}}(t)\,\loss_{\mathrm{geo}},
\end{aligned}
\end{equation}
with \(\lambda_{\mathrm{ord}}=0.1\) fixed, and ramps \(r_{\cdot}(t)\) defined in
Section~\ref{sec:ramps}. The geodesic term itself is
\(\loss_{\mathrm{geo}}=\loss_{\mathrm{anchor}}+0.25\,\loss_{\mathrm{stress}}\).

%==============================================================================
\section{Graph construction}
\label{sec:graph}
%==============================================================================

\subsection{\texorpdfstring{\(\varepsilon\)-net}{epsilon-net} representatives}
\label{sec:enet}

Exact and near-duplicates inflate \(k\)-NN graphs and waste capacity. leanmap
optionally collapses \(X\) to an \(\varepsilon\)-net of representatives before
neighbour search.

\begin{definition}[Estimated \(\varepsilon\)]
\label{def:epsilon}
Let \(r_1(i)\) be the distance from \(x_i\) to its nearest neighbour, taken over
\emph{all} \(N\) points. Set
\[
\varepsilon
=
Q_{0.01}\bigl(\{r_1(i)\}_{i=1}^N\bigr).
\]
If this quantile is \(\le 0\) (exact-duplicate mass dominates), fall back to
the median of the strictly positive \(1\)-NN distances (or \(10^{-12}\) if
none exist).
\end{definition}

\begin{remark}[Why the quantile is taken over all \(N\)]
\label{rem:epsdrift}
Reading the quantile off a fixed-size subsample makes \(\varepsilon\) a
function of dataset size rather than of the data. On an \(m\)-dimensional
manifold the typical \(1\)-NN distance scales as \(n^{-1/m}\), so estimating on
\(n_{\mathrm{sub}}\) of \(N\) points inflates \(\varepsilon\) by
\((N/n_{\mathrm{sub}})^{1/m}\) — a factor of about \(2\) at \(m\approx 6\),
\(N=10^6\), \(n_{\mathrm{sub}}=10^4\), and precisely in the regime where
deduplication is doing real work. The full pass is exact for
\(N\le 2\times 10^4\); above that an ANN index supplies the candidate
neighbour and the distance is still evaluated exactly. When no such index is
available (a metric with no Euclidean embedding, no backend installed) the
implementation retains a subsample but multiplies the quantile by
\((n_{\mathrm{sub}}/N)^{1/m}\), with \(m\) the Levina--Bickel
estimate~\eqref{eq:lb}, which removes the leading size dependence. The path
taken is recorded in the graph statistics.
\end{remark}

The value resolved at fit time is written back into the saved configuration.
A refit therefore reuses the stored radius instead of re-estimating one, which
makes \(\varepsilon\) \emph{absolute in effect after the first fit}: the merge
scale is a fixed property of the artefact, reached through a data-driven
estimator rather than a hand-set tolerance.

Within each top-\(1\) landmark bucket, a greedy \(\varepsilon\)-net is grown
in random order: accept \(x_i\) into an existing cell with representative
\(r_j\) iff \(\min_j d_X(x_i,x_{r_j})\le\varepsilon\); otherwise open a new
cell. Cell weight \(w_i\) equals the number of raw members. Compression ratio
\(N/R\) near \(1\) means the net was effectively a no-op.

\subsubsection*{Halo pass across bucket boundaries}
\label{sec:halo}

Growing the net inside top-\(1\) buckets leaves a gap: a near-duplicate pair
separated by a Voronoi boundary is never compared, so both points survive as
distinct cells. Those are not incidental survivors — the bucket boundary is
exactly where the conditioning code \(a(x)\) switches, so a duplicate that
straddles it is the one most likely to receive two different FiLM
modulations.

A single post-pass closes them. Candidate pairs are drawn from the top-\(c\)
landmark shortlist already computed for the inverted index, restricted to pairs
in \emph{different} top-\(1\) buckets, and accepted when
\(d_X\le\varepsilon\). All candidates are evaluated against the
\emph{pre-merge} representative set, and the accepted pairs are then collapsed
with union--find. The outcome is therefore the connected-component partition of
the accepted-pair graph, which does not depend on the order in which pairs are
visited; a chain \(A\!-\!B\!-\!C\) yields one cell rather than whichever pair
happened to be seen first. Component roots are the lowest representative index,
so the labelling is reproducible across seeds. Cell weights add.

One invariant is deliberately given up: a merged cell no longer lies inside a
single landmark bucket. Every consumer of the bucket-to-cell map — the inverted
index shortlist, cell-to-raw-member expansion in the edge sampler, and ordinal
mid-sampling — either recomputes bucket assignments from the surviving
representatives after the merge, or indexes raw points rather than cells, so
none of them depends on the discarded invariant.

\subsection{Landmarks}
\label{sec:landmarks}

Landmarks \(\{M_\ell\}_{\ell=1}^L\) serve three roles: (i)~conditioning codes
for the encoder, (ii)~support of the cover score, and (iii)~nodes of the
geodesic backbone.

Three selection schemes are available:

\begin{enumerate}[leftmargin=*]
\item \textbf{Ambient farthest-point sampling (FPS).}
  Start from a random index; iteratively
  \(j\leftarrow\arg\max_i\min_{\ell\in\mathrm{chosen}}d_X(x_i,M_\ell)\).
\item \textbf{Geodesic FPS.}
  The same greedy max--min rule on shortest-path distances in a symmetrised
  \(k\)-NN graph of \(X\). Unreachable nodes are treated as farther than any
  finite path so every connected component is covered.
\item \textbf{Geodesic Poisson-disk (blue-noise).}
  Estimate a packing radius \(r\) from a short geodesic-FPS pass of \(L\)
  points (coverage radius \(\times\) a scale), then accept points in random
  order iff their geodesic distance to every accepted point is at least \(r\).
  Unlike FPS, this yields an interior-uniform Delone set with a guaranteed
  minimum geodesic separation.
\end{enumerate}

Each training point is assigned to its \(c=\texttt{C\_BUCKETS}=8\) nearest
landmarks for IVF shortlists and ordinal mid-sampling.

\subsection{\texorpdfstring{\(k\)-nearest}{k-nearest} neighbours}
\label{sec:knn}

On the \(R\) representatives, a \(k\)-NN graph is built in one of
\texttt{brute} / \texttt{ivf} / \texttt{ann} modes (\texttt{auto} chooses by
scale). IVF candidates are the union of the buckets of the query's top
\(c_{\mathrm{search}}=8\) landmarks; ANN uses FAISS IVF-Flat or HNSW with
over-fetch and exact re-rank. Recall against a brute probe set is checked
and the search budget is increased if recall \(<0.9\).

\subsection{\texorpdfstring{Smooth \(k\)-NN}{Smooth k-NN} and fuzzy memberships}
\label{sec:smoothknn}

Following UMAP's fuzzy simplicial set construction~\citep{mcinnes2018umap}:

\begin{definition}[Local connectivity and bandwidth]
For each representative \(i\), let \(\rho_i\) be the
\texttt{local\_connectivity}-th smallest \emph{strictly positive} \(k\)-NN
distance. Solve \(\sigma_i>0\) by bisection so that
\begin{equation}
\label{eq:sigma}
\sum_{j\in\mathcal{N}_k(i)}
\exp\!\Bigl(-\frac{\max(d_{ij}-\rho_i,0)}{\sigma_i}\Bigr)
=
\log_2 k.
\end{equation}
A floor \(\sigma_i\ge\texttt{min\_sigma\_frac}\cdot\overline{\rho}\) prevents
degenerate rows.
\end{definition}

Directed memberships are then
\begin{equation}
\label{eq:pij}
p_{ij}
=
\exp\!\Bigl(-\frac{\max(d_{ij}-\rho_i,0)}{\sigma_i}\Bigr),
\qquad j\in\mathcal{N}_k(i).
\end{equation}

\subsection{Fuzzy union, multiplicity, and backbone}
\label{sec:fuzzyunion}

Let \(P\) be the sparse directed matrix of \(p_{ij}\). Symmetrisation uses the
probabilistic \(t\)-conorm (fuzzy union), \emph{not} mutual \(k\)-NN:
\begin{equation}
\label{eq:fuzzyunion}
P^{\mathrm{sym}}
=
P + P^\top - P\circ P^\top.
\end{equation}
Mutual \(k\)-NN orphans sparse regions; the fuzzy union retains one-sided
memberships.

Cell multiplicities re-enter as
\begin{equation}
\label{eq:multiplicity}
P^{\mathrm{sym}}_{ij}
\;\leftarrow\;
P^{\mathrm{sym}}_{ij}\,(w_i w_j)^{\beta},
\qquad
\beta=\texttt{beta\_multiplicity},
\end{equation}
followed by normalisation \(\max_{ij}P^{\mathrm{sym}}=1\).

\begin{remark}[\(\beta\) is a modelling choice, not plumbing]
\label{rem:beta}
The exponent \(\beta\) states what a repeated row \emph{means}, which is a
property of the dataset and not of the solver, so it is a configuration field
rather than a module constant. At \(\beta\to 0\) multiplicity is ignored and an
\(\varepsilon\)-cell contributes as a single point: the right choice when
duplicates are a deposition artefact carrying no density information, such as
repeated PDB submissions of the same structure or redundant crystal forms of
one compound. At \(\beta\to 1\) the reweighting is proportional and
multiplicity is treated as genuine sampling density that the layout should
reproduce. The default \(\beta=0.5\) splits the difference, letting a cell of
weight \(w\) act like \(\sqrt{w}\) points. Note the interaction with
\(\varepsilon\): the two together decide how much density a duplicate-heavy
region is credited with, so a change to one should be read against the other.
\end{remark}

An optional landmark MST backbone inserts edges of fixed weight
\(\lambda_{\mathrm{backbone}}=0.01\) between landmarks (mapped to nearest
representatives), merged into \(P^{\mathrm{sym}}\) by elementwise maximum and
renormalised. The training edge list retains the upper triangle
(\(i<j\), weight \(>0\)).

\subsection{Multi-scale Galerkin pyramid}
\label{sec:pyramid}

Coarse levels supply long-range attraction so distant regions do not drift
apart.

\begin{definition}[Level sizes]
With finest representative count \(R_0\) and ratio
\(\rho=\texttt{PYRAMID\_REP\_RATIO}=4\),
\[
R_\ell^{\mathrm{target}}
=
\max\Bigl(\mathrm{round}(R_0/\rho^\ell),\;
\texttt{pyramid\_min\_reps}\Bigr),
\quad \ell=1,\ldots,\texttt{pyramid\_scales}.
\]
Coarsening stops when the target is no smaller than the previous level.
\end{definition}

Coarsening of level \(\ell\) to \(\ell{+}1\):

\begin{enumerate}[leftmargin=*]
\item Coarse representatives = FPS among current representatives.
\item Every fine representative is assigned to its nearest coarse
  representative; CSR maps compose down to raw points.
\item \textbf{Galerkin aggregation.} A coarse edge weight is the sum of
  fine-edge weights that cross between the two coarse cells; intra-cell
  (self-loop) edges are dropped. Connectivity is therefore inherited from
  fine edges---no Isomap shortcuts across folds are invented.
\item Aggregated weights are heavy-tailed and must be squashed into a
  membership in \((0,1)\). The default (\texttt{pyramid\_squash} \(=\)
  \texttt{rational\_q99}) is the monotone rational map anchored at the
  \(0.99\)-quantile \(q_{99}\),
  \begin{equation}
  \label{eq:squash}
  \tilde w = \frac{w/q_{99}}{w/q_{99}+c},
  \qquad c=\tfrac{1}{99},
  \end{equation}
  so that \(w=q_{99}\mapsto 0.99\). The plain variant \(w/(w+q_{50})\) and the
  legacy quantile-clamp are both selectable.
\end{enumerate}

If the coarsest level is disconnected, exactly \(n_{\mathrm{comp}}-1\) bridge
edges of weight \texttt{pyramid\_\allowbreak coarse\_\allowbreak backbone} are
added via a union-find seeded with existing components (candidates from an MST
over representative distances). This joins islands without overwriting existing
edge weights.

During training the per-step edge budget \(B_e=\texttt{batch\_edges}\) is
split across levels proportionally to
\texttt{pyramid\_level\_weights} \(=(w_0,\ldots,w_{L_{\mathrm{pyr}}-1})\)
(finest first). Sampled edges from level \(\ell\) are further multiplied by
\(w_\ell\). Coarse-heavy schedules (e.g.\ \((1,2,8)\)) improve geodesic and
density Spearman relative to flat weights.

\subsection{Choosing the squash}
\label{sec:squash}

\begin{remark}[Why not quantile-normalise and clamp]
\label{rem:squash}
The original rule divided by \(q_{99}\) and clamped to \([0,1]\).
Max-normalisation is indeed wrong --- it would squash the bulk toward zero ---
but the clamp variant flattens \emph{every} coarse edge above the \(0.99\)
quantile to a common weight of exactly \(1\). Those are precisely the strongest
long-range bridges, the edges a coarse-heavy schedule such as \((1,2,8)\)
exists to exploit; the rule discarded the ordering among them. Any rational
map is strictly monotone and never saturates, so the ranking survives intact.
\end{remark}

Monotonicity alone was not sufficient, and the measurement is worth recording
because the obvious fix regressed badly. Replacing the clamp with the
median-anchored \(w/(w+q_{50})\) cost \(0.21\) of digits density Spearman ---
five times the pooled seed spread, and by far the largest resolved effect
anywhere in this note. Density structure is exactly what the pyramid is
supposed to supply, so the strictly-better-shaped map made the map
substantially worse at its own job.

The re-tune anticipated above was run and did \emph{not} rescue it. Sweeping
\texttt{pyramid\_level\_weights} across flat \((1,1,1)\), matched, ramp
\((1,2,8)\), and steep schedules moves density Spearman between \(0.27\) and
\(0.56\), all resolved-worse than the clamp's \(0.725\). So the loss is not a
level-weight mismatch that a compensating schedule can absorb.

The cause is where in the distribution the map is anchored. Sending the
\emph{median} coarse edge to \(0.5\) lifts the bulk of the coarse graph to a
strength it never had under the clamp, and the strong bridges --- which under
the clamp stood out by magnitude as well as rank --- are compressed into the
same range as everything else. Anchoring the same rational map at \(q_{99}\)
keeps the bulk small and sends the top percentile to \(0.99\) without ever
reaching \(1\), preserving both the magnitude profile and strict monotonicity:

\begin{center}
\small
\begin{tabular}{lcccc}
\toprule
Squash & density Spearman & geodesic Sp. & spacing CV & monotone? \\
\midrule
\texttt{quantile\_clamp}       & \(0.725\pm0.009\) & \(0.712\pm0.015\) & \(0.583\pm0.084\) & no \\
\texttt{rational} (\(q_{50}\)) & \(0.520\pm0.039\) & \(0.724\pm0.058\) & \(0.389\pm0.034\) & yes \\
\texttt{rational\_q99}         & \(0.678\pm0.017\) & \(0.689\pm0.075\) & \(0.477\pm0.057\) & yes \\
\bottomrule
\end{tabular}
\end{center}

\(q_{99}\)-anchoring recovers about \(77\%\) of the lost density fidelity. The
remaining \(-0.047\) against the clamp is still resolved, and that is the price
of monotonicity: some of what the clamp bought was the flattening itself, not
merely the magnitude profile. It is a price worth paying, since the clamp's
advantage comes from discarding the ordering among the strongest long-range
bridges, and it comes with a compensating resolved gain --- spacing CV falls
from \(0.583\) to \(0.477\). Geodesic Spearman, label accuracy,
trustworthiness, and \(k\)-NN overlap are all unresolved between the three.

\texttt{rational\_q99} is therefore the default, and the level weights
\((1,2,8)\) transfer unchanged --- which was the point of matching the
magnitude profile rather than re-tuning around a mismatched one.

%==============================================================================
\section{Parametric model}
\label{sec:model}
%==============================================================================

\subsection{Anchor affinity}
\label{sec:affinity}

\begin{definition}[Softmax landmark affinity]
With learnable temperatures \(\tau_\ell=\tau_{\min}+e^{\log\tau_\ell}\)
(\(\tau_{\min}=10^{-3}\)),
\begin{equation}
\label{eq:affinity}
D_{i\ell}=d_X(x_i,M_\ell),
\qquad
a_{i\ell}
=
\frac{\exp(-D_{i\ell}/\tau_\ell)}
{\sum_{\ell'}\exp(-D_{i\ell'}/\tau_{\ell'})}.
\end{equation}
\end{definition}

Default \(\tau_\ell\) equals the mean distance from landmark \(\ell\) to its
\(\min(32,L-1)\) nearest landmarks, times \texttt{tau\_scale}. Anchors are
frozen automatically if \(d_X\) is non-differentiable in its second argument.

Affinity perplexity
\[
\mathrm{perp}_i
=
\exp\Bigl(-\sum_\ell a_{i\ell}\log a_{i\ell}\Bigr)
\]
is monotone in \(\tau\); \texttt{calibrate.py} bisects \texttt{tau\_scale} so
that the median perplexity meets a target (default \(8\)). Perplexity near
\(1\) is a Voronoi indicator; near \(L\) the code is uninformative.

\subsection{FiLM encoder}
\label{sec:film}

Let \(x_n=(x-\mu)/\sigma\) be the encoder-normalised input. A depth-\(K\)
MLP with width \(W\) applies, at each layer \(k=1,\ldots,K\),
\begin{equation}
\label{eq:film}
\begin{aligned}
h_0 &= x_n,\\
\tilde h_k &= W_k h_{k-1}
  \qquad\text{(bias-free)},\\
h_k &= \GELU\Bigl(\gamma_k\odot \LN\bigl(\tilde h_k\bigr) + \beta_k\Bigr),
\end{aligned}
\end{equation}
with FiLM parameters \((\gamma_k,\beta_k)\in\R^W\) produced from the affinity
code (Section~\ref{sec:factors}). The embedding is
\begin{equation}
\label{eq:z}
z
=
\begin{cases}
W_{\mathrm{PCA}}\,x_n + W_{\mathrm{head}}\,h_K + b
  & \text{if \texttt{pca\_skip}=True},\\[4pt]
W_{\mathrm{head}}\,h_K + b
  & \text{otherwise}.
\end{cases}
\end{equation}
When the PCA skip is enabled, \(W_{\mathrm{PCA}}\) holds the top-\(d\) right
singular vectors of the (optionally centred) normalised training matrix, and
the residual head is initialised \(\mathcal{N}(0,10^{-4})\) with zero bias, so
training \emph{starts at plain PCA}.

\begin{remark}[Coupled \texttt{pca\_skip} / learning rate]
With the skip on and \(\mathrm{lr}=10^{-3}\), the unconstrained linear path
tends to dominate the layout. Turning the skip off at the same learning rate
undertrains the near-zero head. Escaping both regimes requires
\texttt{pca\_skip}=\texttt{False} \emph{and} a raised learning rate
(\(\sim 2\cdot 10^{-2}\) for \(N\le 5\cdot 10^3\)). The two knobs are one
decision~\cite{leanmap-config}.
\end{remark}

\subsection{Factor roles and FiLM composition}
\label{sec:factors}

Multiple conditioning factors may be stacked. Each factor \(f\) emits a
contribution \((\gamma^{(f)},\beta^{(f)})\) according to a role:

\begin{center}
\small
\begin{tabular}{lll}
\toprule
Role & \(\gamma^{(f)}\) & \(\beta^{(f)}\) \\
\midrule
\texttt{PRIMARY} & per-channel \((B,K,W)\) & per-channel \\
\texttt{MODULATOR} & per-layer scalar & per-channel \\
\texttt{GAIN} & per-layer scalar & --- \\
\texttt{AXIS} & --- (monotone output axis) & --- \\
\bottomrule
\end{tabular}
\end{center}

Composition is multiplicative in scale and additive in shift,
\begin{equation}
\label{eq:filmcompose}
\gamma
=
\clamp\Bigl(\prod_f \gamma^{(f)},\;\gamma_{\min},\;\gamma_{\max}\Bigr),
\qquad
\beta
=
\sum_f \beta^{(f)},
\end{equation}
with \(\gamma_{\min}=0.1\), \(\gamma_{\max}=10\). The clamp hit-rate is
monitored as an extrapolation diagnostic.

\begin{remark}[Why FiLM must follow the normalisation]
\label{rem:filmorder}
Equation~\eqref{eq:film} normalises \emph{before} modulating. The opposite
order, \(\GELU(\LN(\gamma\odot\tilde h+\beta))\), silently disables two of the
roles: layer normalisation satisfies \(\LN(c\,h)=\LN(h)\) for any \(c>0\), so a
per-layer \emph{scalar} \(\gamma\) is exactly invisible. A \texttt{GAIN} factor
would then contribute nothing to the forward pass while still receiving
gradient noise, and a \texttt{MODULATOR} would act only through the ratio
\(|\gamma\,\tilde h|:|\beta|\) rather than through its scale at all.
\texttt{PRIMARY} is unaffected either way, because a per-channel rescale
changes the cross-channel pattern that \(\LN\) normalises. Earlier revisions of
this package used the modulate-first order; role comparisons involving
\texttt{GAIN} or \texttt{MODULATOR} predate the fix and are not comparable
across it, whereas single-\texttt{PRIMARY} results (the default path) remain
qualitatively valid.
\end{remark}

An \texttt{AXIS} factor contributes a monotone scalar MLP (non-negative
weights via \(\softplus\)) added to a dedicated output coordinate.

%==============================================================================
\section{Low-dimensional similarity and \texorpdfstring{\texttt{min\_dist}}{min\_dist}}
\label{sec:ab}
%==============================================================================

\subsection{Curve fit}
\label{sec:curvefit}

As in UMAP, the low-dimensional similarity kernel is parameterised by
\((a,b)\) fitted to a piecewise target. On a grid \(x\in[0,3\cdot\mathrm{spread}]\),
\begin{equation}
\label{eq:target}
y(x)
=
\begin{cases}
1, & x < \texttt{min\_dist},\\
\exp\bigl(-(x-\texttt{min\_dist})/\mathrm{spread}\bigr),
  & x\ge\texttt{min\_dist},
\end{cases}
\end{equation}
and
\begin{equation}
\label{eq:curve}
\hat y(x)
=
\frac{1}{1+a\,x^{2b}}.
\end{equation}
Parameters are obtained by bounded least squares
(\(a,b>0\), initialisation \((1,1)\)). Fallback:
\((a,b)=(1.577,0.895)\) (UMAP's \texttt{min\_dist}\,{=}\,0.1 values).
Default \(\mathrm{spread}=1\).

The embedding kernel used in the geometric loss is then
\begin{equation}
\label{eq:q}
q(u,v)
=
\bigl(1 + a\,\norm{u-v}^{2b}\bigr)^{-1},
\end{equation}
clamped to \([10^{-7},1-10^{-7}]\) for logarithms.

\subsection{Stability of the attractive force}
\label{sec:bstability}

\begin{proposition}[Near-contact attraction]
For small separation \(r=\norm{z_i-z_j}\), the attractive force derived from
\(-\log q\) scales as \(2ab\,r^{2b-1}\). The force \emph{per unit separation}
is therefore \(\propto r^{2b-2}\), which diverges as \(r\to 0\) whenever
\(b<1\): a pair that is already close is pulled proportionally harder.
\end{proposition}

This positive feedback is the mechanism that collapses neighbourhoods into
knots separated by voids, and it is a statement about the ratio of force to
separation, not about the force itself.

\begin{remark}[Correcting the ``hard floor'' argument]
\label{rem:bfloor}
An earlier version of this note argued that because leanmap differentiates the
loss directly rather than using UMAP's clipped SGD kernel, \(b\ge 1\) is a hard
floor. That mechanism does not hold. The attractive force \(2ab\,r^{2b-1}\)
diverges at contact only for \(b<\tfrac12\); for \(\tfrac12<b<1\) --- which is
the entire range in question, since UMAP's default \(b\approx 0.895\) sits
there --- it \emph{vanishes} sublinearly. There is nothing for autodiff to blow
up on.

What actually diverges in UMAP is the \emph{repulsive} gradient, whose
\(1/(\varepsilon+r^2)\) factor grows without bound as a negative sample
approaches its anchor, and that is what UMAP's clipping addresses. It is a
property of the repulsive term, not a consequence of clipped versus exact
differentiation of the attractive one.

The consequence is that \(b\ge 1\) has no derivation behind it. What survives
is the proposition above --- sub-unit \(b\) does produce a relative pull that
grows without bound at contact --- plus whatever the measurements support. Any
threshold on \(b\) is stated below as an empirical rule.
\end{remark}

\paragraph{The seeded ladder.}
An earlier version of this table was single-seed, and reported digits 5-NN of
\(0.958/0.969/0.947/0.939\) across \(\texttt{min\_dist}=0.1/0.2/0.5/0.8\) --- a
spread of \(0.03\) --- from which it concluded that accuracy ``degrades
monotonically past \(0.2\)''. Re-run at three seeds, with the seed spread
reported (\(\pm\) is the sample sd over seeds, \(n=3\)):

\begin{center}
\small
\begin{tabular}{cc|ccc|cc}
\toprule
& & \multicolumn{3}{c|}{digits} & \multicolumn{2}{c}{s-curve} \\
\texttt{min\_dist} & \(b\) & 5-NN & trust\(_{15}\) & spacing CV
  & \(k\)NN overlap & area sd \\
\midrule
0.1 & 0.895 & \(0.936\pm0.029\) & \(0.948\pm0.009\) & \(0.612\pm0.110\)
    & \(0.813\pm0.006\) & \(0.231\pm0.022\) \\
0.2 & 1.003 & \(0.929\pm0.028\) & \(0.944\pm0.009\) & \(0.504\pm0.077\)
    & \(0.824\pm0.019\) & \(0.193\pm0.019\) \\
0.5 & 1.334 & \(0.921\pm0.034\) & \(0.938\pm0.010\) & \(0.398\pm0.050\)
    & \(0.858\pm0.016\) & \(0.152\pm0.017\) \\
0.8 & 1.681 & \(0.906\pm0.030\) & \(0.921\pm0.011\) & \(0.334\pm0.028\)
    & \(0.858\pm0.041\) & \(0.141\pm0.032\) \\
\bottomrule
\end{tabular}
\end{center}

\begin{remark}[The accuracy trend is real but unresolved; the uniformity trend
is large]
\label{rem:mindistseeded}
Digits 5-NN does decline monotonically, by \(0.030\) across the full ladder ---
but the seed sd is itself \(0.028\)--\(0.034\), so not one pairwise difference
in that column is resolved, including the endpoints. The direction is
consistent with the single-seed table; the magnitude is not established. What
the earlier text called ``degrades monotonically past \(0.2\)'' should be read
as a trend the data cannot presently distinguish from noise. Note also that the
seeded means sit \emph{below} the single-seed numbers at every rung, which is
what one expects when a single run is implicitly the best of an unrecorded few.

The uniformity trend, by contrast, is large and resolved: digits spacing CV
falls from \(0.612\) to \(0.334\), with \(0.5\) and \(0.8\) both resolved
against \(0.1\). On the s-curve, \(\texttt{min\_dist}=0.5\) is resolved
\emph{better} than \(0.1\) on \(k\)-NN overlap (\(+0.044\)),
trustworthiness, and \(\mathrm{area\_sd}\) (\(-0.079\)) --- on a smooth
manifold, raising \texttt{min\_dist} improves neighbourhood preservation
outright rather than trading it away.
\end{remark}

The default \(\texttt{min\_dist}=0.5\) is retained, and the ladder now says why
rather than asserting a regime boundary. It is the largest value with no
resolved loss anywhere: at \(0.8\), digits trustworthiness is resolved-worse
(\(-0.027\)), while the additional uniformity over \(0.5\) is small and the
s-curve metrics stop improving. Going the other way, \(0.1\) and \(0.2\) give
up the whole uniformity gain and, on the s-curve, real neighbourhood fidelity,
for an accuracy difference the seeds cannot resolve.

Two claims from the earlier text are withdrawn: ``values below \(0.2\) are
never safe'', which the \(0.1\) row does not support, and the collapse /
marginal / stable regimes stated as thresholds on \(b\), which rested on the
mechanism corrected in Remark~\ref{rem:bfloor}. The surviving rule is
empirical: prefer the top of the ladder, back off if trustworthiness degrades.

%==============================================================================
\section{Training objective}
\label{sec:loss}
%==============================================================================

\subsection{Fuzzy cross-entropy}
\label{sec:geom}

Positive edges \((i,j)\) are drawn proportional to \(P^{\mathrm{sym}}_{ij}\)
(Vose alias sampling), expanded from \(\varepsilon\)-net cells to random raw
members. Negatives \(z_i^-\) are uniform over representatives
(\(n_{\mathrm{neg}}=5\) by default).

\begin{equation}
\label{eq:geom}
\begin{aligned}
\loss_{\mathrm{geom}}
&=
-\frac{1}{B}\sum_{(i,j)\in E_B} w_{ij}\log q(z_i,z_j) \\
&\quad
-\frac{1}{B\,n_{\mathrm{neg}}}
\sum_i\sum_{m=1}^{n_{\mathrm{neg}}}
\log\bigl(1-q(z_i,z_{i,m}^-)\bigr).
\end{aligned}
\end{equation}

This is the UMAP-style cross-entropy between fuzzy memberships and the
low-dimensional kernel~\eqref{eq:q}, with direct autodiff rather than a
clipped-force approximation.

\subsection{Ordinal triplet loss}
\label{sec:ord}

For verified quadruples (anchor, near, mid, far) with
\(d_{\mathrm{an}}<d_{\mathrm{am}}<d_{\mathrm{af}}\) in a factor's view
metric, with adaptive scale \(s\leftarrow 0.9\,s+0.1\,\overline{d_{\mathrm{af}}}\),
\begin{equation}
\label{eq:ord}
\loss_{\mathrm{ord}}
=
\E\Bigl[
-\log\sigma\!\Bigl(\frac{d_{\mathrm{am}}-d_{\mathrm{an}}}{s}\Bigr)
-\log\sigma\!\Bigl(\frac{d_{\mathrm{af}}-d_{\mathrm{am}}}{s}\Bigr)
\Bigr].
\end{equation}
Mid points are drawn from the bucket of the anchor's rank-\(r\) landmark with
log-uniform \(r\sim\mathrm{round}(\exp\mathcal{U}(\log 2,\log L))\). The
weight \(\lambda_{\mathrm{ord}}=0.1\) is fixed plumbing.

\paragraph{Retention and its null.}
Retention is the fraction of proposed triplets that survive verification. It is
interpretable only against the rate at which a triplet would pass with
\emph{no} information in the landmark ranking, and that rate is not a universal
constant: it depends on the sampler, on the bucket-size distribution, and on
the data. The note previously asserted a chance level of \(0.475\), which is
not recoverable from the log-uniform scheme above.

The null is therefore \emph{measured} at fit time by permuting each row's
landmark ordering and re-running the identical sampler, which destroys the
rank--distance association while leaving the rank draw, the bucket-based mid
selection, and the uniform far point untouched. Measured values are \(0.472\)
(iris), \(0.484\) (swiss roll), and \(0.487\) (digits) --- close to the old
constant, but data-dependent, and correct by construction if the sampling
scheme ever changes. The warning threshold floats with the measured null at the
same margin. That margin is narrow: a retention of \(0.55\) sits only
\(\approx 0.07\) above chance, so values in the \(0.5\)--\(0.6\) band should be
read as barely informative rather than as a pass.

\subsection{Landmark regularisation}
\label{sec:lm}

\begin{equation}
\label{eq:lm}
\loss_{\mathrm{lm}}
=
\frac{1}{B}\sum_i\sum_\ell a_{i\ell}\,D_{i\ell}
-
\eta\sum_\ell \bar a_\ell\log\bar a_\ell,
\qquad
\bar a_\ell=\frac{1}{B}\sum_i a_{i\ell},
\end{equation}
with \(\eta=\texttt{ETA\_BALANCE}=1\). The first term is soft quantisation
error; the second is a load-balancing entropy over mean affinities.

\subsection{Local rigidity (frame term)}
\label{sec:frame}

The frame term is an as-rigid-as-possible~\citep{sorkine2007arap} penalty in
the Gram domain. For a centre \(c\) with neighbours \(\{k\}\) (a
\emph{star}), write offsets \(u_k=x_k-x_c\) and \(v_k=z_k-z_c\). A similarity
preserves all pairwise inner products up to one scalar, so
\begin{equation}
\label{eq:grams}
(G_x)_{kl}=\inner{u_k}{u_l},\qquad
(G_z)_{kl}=\inner{v_k}{v_l}.
\end{equation}
With pair mask \(M_{kl}=m_k m_l\) (padding / dropped neighbours),
\begin{equation}
\label{eq:frame}
s^2
=
\frac{\inner{G_z}{G_x}_M}{\inner{G_x}{G_x}_M}
\quad\text{(detached)},
\qquad
\loss_{\mathrm{frame}}
=
\frac{\sum\norm{G_z-s^2 G_x}_M^2}
{\sum\norm{G_z}_M^2}.
\end{equation}
The diagonal recovers edge-length isometry; off-diagonals encode relative
orientation, so a length-preserving shear (frame twist) is penalised. The
objective is rotation/reflection-invariant and needs no SVD in the backward
pass. Scale invariance follows from the detached \(s^2\) and the
normalisation by \(\norm{G_z}^2\).

\paragraph{Tangent / geodesic mode.}
On folded manifolds, ambient \(k\)-NN includes across-sheet shortcuts. With
\texttt{frame\_tangent}=\texttt{True}, under \(\mathrm{no\_grad}\):

\begin{enumerate}[leftmargin=*]
\item Distance-weight offsets
  \(w_k=\exp(-\tfrac12(\norm{u_k}/h)^2)\) with \(h\) the mean neighbour
  distance.
\item Estimate a local tangent basis as the top-\(d_{\mathrm{eff}}\) right
  singular vectors of the weighted offset cloud.
\item Drop neighbours with normal residual fraction
  \(\norm{u-\Pi_T u}/\norm{u}>\texttt{FRAME\_NORMAL\_THRESH}=0.5\).
\item Match Grams on tangent-projected offsets \(u\leftarrow u\,T\).
\end{enumerate}
This yields a first-order geodesic rigidity safe on swiss-roll-like data.

In this mode the term is essentially LTSA~\citep{zhang2004ltsa}: steps 1--2 are
LTSA's weighted local tangent estimate, and step 4 imposes LTSA's requirement
that the tangent-coordinate structure agree between data and embedding. The
departures are that the agreement is enforced on Gram matrices with a detached
scale rather than by aligning coordinates, and that it is a penalty on a
parametric map rather than a global eigenproblem solved for the embedding.

\paragraph{Why the ramp must be delayed on fold-back manifolds.}
Local rigidity cannot distinguish a rolled manifold from its unrolled
isometry: it penalises the \emph{transient} stretch of unrolling. On such
data, \(r_{\mathrm{frame}}\) should rise only after the neighbour loss has
globally unrolled (e.g.\ \((0.5,0.75)\)); it then cleans up width and twist.
On non-folding sheets early or delayed both help; delayed remains safer.

An older length-only local isometry loss (pairwise \(\norm{z_i-z_j}\) vs.\
\(\norm{x_i-x_j}\) up to detached scale) is implemented but unused in the
default trainer; the Gram term supersedes it.

\subsection{Geodesic backbone}
\label{sec:geo}

Landmarks are mapped to nearest training rows; a symmetrised \(k\)-NN graph
on \(X\) yields landmark--landmark shortest paths \(g_{\ell\ell'}\). Classical
MDS~\citep{torgerson1952,isomap2000} produces a frozen target \(Y\in\R^{L\times d}\):
\begin{equation}
\label{eq:cmds}
J=I-\tfrac1L\ind\ind^\top,\qquad
B=-\tfrac12 J\,G^{(2)}\,J,\qquad
Y=V_d\,\Lambda_d^{1/2},
\end{equation}
with \(\Lambda_d=\max(\lambda_d,0)\) from the eigendecomposition of the
symmetrised Gram \(B\).

\begin{remark}[The clamp discards evidence; report it]
\label{rem:negeigen}
\(B\) is positive semi-definite only when the geodesic distances \(g\) are
Euclidean-embeddable. They generally are not --- a manifold with a hole, or one
whose shortest paths wrap around a fold, produces genuinely non-Euclidean
\(g\), and the negative eigenvalues of \(B\) are precisely the record of that.
The clamp \(\max(\lambda,0)\) throws the record away silently, leaving a target
\(Y\) that looks unremarkable no matter how badly it fits.

The negative mass is therefore emitted as a first-class diagnostic,
\begin{equation}
\label{eq:negeigen}
\nu = \frac{\sum_{\lambda_i<0}\abs{\lambda_i}}{\sum_i\abs{\lambda_i}},
\end{equation}
alongside the fraction of positive spectrum captured by the retained \(d\)
components. Large \(\nu\) means the MDS target is a poor summary of the
geodesic structure and \(\loss_{\mathrm{anchor}}\) is pulling towards
something the data does not support.

The response is a warning, not an adjustment. Making
\(\lambda_{\mathrm{anchor}}\) a hidden function of \(\nu\) would reintroduce
exactly the buried-coupling problem that motivated exposing \(\beta\)
(Remark~\ref{rem:beta}): the recipe would no longer say what the objective is.
The diagnostic is reported, a threshold triggers a warning, and the weight
stays under the user's control.
\end{remark}

\paragraph{Scale-invariant stress.}
On sampled landmark pairs,
\begin{equation}
\label{eq:stress}
s=\frac{\inner{d_z}{g}}{\inner{g}{g}}
\quad\text{(detached)},
\qquad
\loss_{\mathrm{stress}}
=
\frac{\E[(d_z-s\,g)^2]}{\E[d_z^2]},
\end{equation}
where \(d_z=\norm{z_a-z_b}\). This constrains \emph{how far} landmarks sit---the
global metric gauge that affinity losses leave free (the ``banana'' bend of an
otherwise isometric unrolling).

\paragraph{Procrustes anchor.}
Let \(Z_M\in\R^{L\times d}\) be current landmark embeddings. Under
\(\mathrm{no\_grad}\), optimally align the frozen MDS target \(Y\) to \(Z_M\)
by a similarity:
\begin{equation}
\label{eq:procrustes}
\begin{aligned}
M &= Y_c^\top Z_{M,c},
\quad
U\Sigma V^\top=\mathrm{SVD}(M),
\quad
R=UV^\top
\quad\text{(reflection-fixed if \(\det R<0\))},\\
c &= \frac{\mathrm{tr}\,\Sigma}{\norm{Y_c}_F^2},
\quad
t = \bar z_M - c\,\bar Y R,\\
\loss_{\mathrm{anchor}}
&=
\frac{1}{Ld}\norm{Z_M - (c\,YR + t)}_F^2.
\end{aligned}
\end{equation}
Gradients flow only into \(Z_M\). Stress alone is invariant to local frame
twists that approximately preserve pairwise distances; the absolute-position
prior selects the untwisted developable layout.

Combined:
\begin{equation}
\label{eq:lgeocombined}
\loss_{\mathrm{geo}}
=
\lambda_{\mathrm{anchor}}\,\loss_{\mathrm{anchor}}
+ 0.25\,\loss_{\mathrm{stress}},
\qquad
\lambda_{\mathrm{anchor}}=1\ \text{by default}.
\end{equation}
The anchor weight is a config knob rather than the hardcoded \(1\) of earlier
versions. Setting \(\lambda_{\mathrm{anchor}}=0\) keeps the scale-invariant
stress while dropping the absolute-position prior, which is worth testing
whenever \(\loss_{\mathrm{frame}}\) is active --- the frame term already
penalises local twist, and does so without committing to a global MDS target
that Remark~\ref{rem:negeigen} may show to be unsupported. The measured
consequence is in \S\ref{sec:ablations}: on swiss roll the means are
indistinguishable, but the seed spread of \(\mathrm{area\_sd}\) grows by
roughly an order of magnitude, so the default stays at \(1\).

\begin{remark}[A Lipschitz penalty was removed]
\label{rem:nolip}
Earlier versions carried a one-sided gradient penalty
\(\E[(\norm{\partial z/\partial\tilde x}_F-1)_+^2]\) on random convex mixtures,
with weight \(\lambda_{\mathrm{lip}}\). It has been deleted rather than
retuned. It was never swept, defaulted to zero in the small-\(N\) recipe that
ships, and — more importantly — it penalises the \emph{absolute} Frobenius norm
of the Jacobian against a hard threshold of \(1\). That contradicts the rest of
the objective: both \(\loss_{\mathrm{frame}}\) and \(\loss_{\mathrm{stress}}\)
are deliberately scale-free, fitting and detaching a global scale \(s\) so that
no term can dictate the overall size of the embedding. A penalty anchored to a
fixed Jacobian norm silently does exactly that. No scale-free replacement was
substituted, because there is no measurement showing the term was needed.
\end{remark}

\subsection{Ramp schedules}
\label{sec:ramps}

For training fraction \(t\in[0,1]\) and breakpoints \((t_0,t_1)\),
\begin{equation}
\label{eq:ramp}
r_\uparrow(t)
=
\begin{cases}
1, & t_0=t_1=0,\\
0, & t<t_0,\\
\dfrac{t-t_0}{t_1-t_0}, & t_0\le t<t_1,\\
1, & t\ge t_1.
\end{cases}
\end{equation}
Default \(\texttt{geo\_ramp}=(0.2,0.45)\); \(\texttt{frame\_ramp}=(0,0)\)
(on immediately) unless delayed for fold-back data.

%==============================================================================
\section{Optimisation}
\label{sec:opt}
%==============================================================================

AdamW with weight decay \(10^{-4}\). Default schedule: linear warmup over
\(5\%\) of total steps (start factor \(0.01\)), then cosine annealing to the
end. Gradients are clipped to \(\ell_2\)-norm \(5\). Non-finite losses skip
the step. Steps per epoch equal \(\lceil E_{\mathrm{fine}}/B_e\rceil\) where
\(E_{\mathrm{fine}}\) is the finest-level edge count.

An optional two-phase override holds a constant learning rate for
\texttt{lr\_switch\_epochs} then switches to \texttt{lr\_after}, disabling
warmup{+}cosine.

%==============================================================================
\section{Evaluation statistics}
\label{sec:eval}
%==============================================================================

\subsection{Shepard / stress}
\label{sec:shepard}

For paired distances \((g_i,e_i)\) (ambient or geodesic vs.\ embedding),
\begin{equation}
\label{eq:shepard}
\alpha=\frac{\inner{g}{e}}{\inner{e}{e}},
\qquad
\mathrm{stress}
=
\sqrt{\frac{\sum_i(\alpha e_i-g_i)^2}{\sum_i g_i^2}},
\end{equation}
plus Spearman rank correlation. Ambient Shepard uses random distinct pairs;
geodesic Shepard uses Dijkstra on the fuzzy graph with hop length
\(1/\clamp(w,10^{-6},\infty)\).

\subsection{Trustworthiness and continuity}
\label{sec:trust}

With ambient ranks \(r^X_{ij}\) and embedding \(k\)-NN sets \(\mathcal{N}^Z_k\),
\begin{equation}
\label{eq:trust}
T_k
=
1-\frac{2}{nk(2n-3k-1)}
\sum_i\sum_{j\in\mathcal{N}^Z_k(i)\setminus\mathcal{N}^X_k(i)}
\bigl(r^X_{ij}-k\bigr),
\end{equation}
and continuity by exchanging the roles of \(X\) and \(Z\). Overlap
\(\mathrm{knn\_overlap}_k\) is the mean fraction of shared neighbours.

\subsection{Density correspondence}
\label{sec:density}

Local density \(\mathrm{dens}(x_i)=(k^{-1}\sum_{j\in\mathcal{N}_k(i)}d_{ij})^{-1}\).
Log-densities \(\ell_a=\log_{10}\mathrm{dens}_X\), \(\ell_z=\log_{10}\mathrm{dens}_Z\)
are compared by OLS slope, residual, Spearman, and Pearson-on-logs.

\subsection{Uniformity}
\label{sec:uniformity}

With \(k\)-NN radii \(r_X(i),r_Z(i)\) (default \(k=10\)) on an interior mask,
\begin{equation}
\label{eq:uniformity}
\mathrm{spacing\_cv}
=
\frac{\mathrm{sd}\,r_Z}{\overline{r_Z}},
\qquad
\mathrm{area\_sd}
=
\mathrm{sd}\Bigl[\log\frac{r_Z}{r_X}\Bigr].
\end{equation}
Zero \(\mathrm{area\_sd}\) means every neighbourhood is blown up by the same
factor (area-preserving up to global scale). A layout can be perfectly uniform
on the page while badly misrepresenting a non-uniform sample; \(\mathrm{area\_sd}\)
is the honest target for that distinction.

\subsection{Persistence: would you get the same map twice?}
\label{sec:persistence}

Every statistic above measures whether \emph{a} map is good. None measures
whether refitting reproduces it --- which is the property a saved
\texttt{.pt} artefact actually sells, since downstream coordinates, stored
cluster labels, and figures all assume the map is stable.

Two maps \(Z,Z'\) of the same points are compared after refits at (i)
\(\ge 3\) seeds and (ii) \(\ge 3\) training subsamples (80\% draws).

\paragraph{Coordinate disagreement.}
The embedding is defined only up to a similarity, so the comparison fits a
Procrustes similarity \((s,Q,t)\) and reports the median residual
\begin{equation}
\label{eq:procdis}
\mathrm{disagree}(Z,Z')
=
\operatorname{median}_i\ \norm{\,sQz_i+t-z'_i\,}\Big/\operatorname{median}_i\norm{z'_i-\bar z'}.
\end{equation}
The transform is fit on a shared \emph{anchor} set and evaluated on
\emph{held-out} points, so the alignment is not scored on the points it was fit
to.

\paragraph{Neighbourhood rank agreement.}
Spearman correlation between the neighbour rankings the two maps induce on a
shared probe set, plus \(k\)-NN Jaccard overlap.

\begin{remark}[Which number is primary]
\label{rem:persprimary}
Rank agreement is the primary number and coordinate disagreement the secondary
one. Procrustes can absorb a rotation but cannot repair genuine topological
ambiguity --- which arm of a branching structure lands on which side is a
discrete choice, and no similarity transform undoes it. A large coordinate
disagreement together with high rank agreement is therefore a gauge artefact,
not instability, and should be reported as such rather than as a failure.
The converse --- small coordinate disagreement with degraded rank agreement ---
is the genuinely bad case.
\end{remark}

\subsection{Matched null}
\label{sec:null}

A configuration is re-fit on column-shuffled input (joint structure destroyed,
marginals retained). Only the margin of a metric over this null carries
information.

\subsection{Intrinsic dimension (calibration)}
\label{sec:id}

Levina--Bickel MLE~\citep{levina2005}:
\begin{equation}
\label{eq:lb}
\hat m
=
\Biggl(\frac1n\sum_i\frac1{k-1}\sum_{j=1}^{k-1}
\log\frac{r_k(i)}{r_j(i)}\Biggr)^{-1}.
\end{equation}

%==============================================================================
\section{Conformal scores and out-of-distribution detection}
\label{sec:conformal}
%==============================================================================

\subsection{Cover score}
\label{sec:cover}

\begin{definition}[Cover]
\begin{equation}
\label{eq:cover}
\mathrm{cover}(x)
=
\min_\ell \frac{d_X(x,M_\ell)}{s_{\mathrm{nat}}},
\end{equation}
where \(s_{\mathrm{nat}}\) is the metric natural scale.
\end{definition}

This is the baseline nonconformity score. A secondary chart-quality diagnostic
compares ambient affinity \(a(x)\) to an embedding-space softmax over landmark
embeddings; it is \emph{not} used as an OOD gate.

\begin{remark}[What may be estimated on all of \(X\), and why]
\label{rem:validity}
The conformal guarantee needs the score function fixed before calibration and
the calibration scores exchangeable with the test scores. Two different
arguments apply here, and conflating them invites a spurious leak report.

\(s_{\mathrm{nat}}\), and any global \((\mu,\sigma)\) standardisation, are
\emph{strictly monotone rescalings} applied identically to calibration and test
scores. Equation~\eqref{eq:pvalue} depends on the scores only through the
count \(\#\{s_j\ge s(x)\}\), which is invariant under any common increasing
map. They may therefore be estimated on all of \(X\) with no effect whatsoever
on the \(p\)-values.

The per-landmark radii and charts of \S\ref{sec:localscale} are \emph{not}
rank-preserving --- they deliberately reorder points --- so the same argument
does not cover them. They are instead fit on the \emph{training} split, which
is disjoint from calibration by construction (\S\ref{sec:overview}: the
calibration split is taken before landmark selection). Both are sound; only the
second needs the split.
\end{remark}

\subsection{Per-landmark local scale}
\label{sec:localscale}

A single global \(s_{\mathrm{nat}}\) makes~\eqref{eq:cover} a statement about
absolute distance, so the sparse tail of the training distribution is flagged
as readily as anything genuinely off-manifold. Replacing the global scale by a
local one,
\begin{equation}
\label{eq:covlocal}
\mathrm{cover}(x)=\min_\ell \frac{d_X(x,M_\ell)}{r_\ell},
\qquad
r_\ell=\operatorname{median}\{\,d_X(x_i,M_\ell): \ell(x_i)=\ell,\ x_i\in\text{train}\,\},
\end{equation}
buys approximate \emph{conditional} validity while keeping one pooled
calibration set: a typical point scores \(\approx 1\) in a dense or a sparse
region alike.

\begin{remark}[Why not Mondrian]
\label{rem:mondrian}
The textbook route to conditional validity is Mondrian conformal prediction,
which pools \emph{calibration} points per bucket. That is not viable at these
shapes: with \(L\sim 64\)--\(256\) landmarks against
\(n_{\mathrm{cal}}\lesssim 200\)--\(2000\), most buckets receive zero or one
calibration point, and a per-bucket \(p\)-value has resolution
\(1/(n_\ell+1)\). The radii in~\eqref{eq:covlocal} and the charts below are fit
from \emph{training} points, which are plentiful, so they do not hit the same
sparsity wall. The asymmetry is deliberate: condition the \emph{score} on the
bucket, keep the \emph{calibration} pooled.
\end{remark}

\subsection{Tangent charts as the default nonconformity}
\label{sec:chartscore}

Isotropic balls are the wrong shape in \(D\gg m\). The acceptance region
\(\{\mathrm{cover}\le\tau\}\) is a union of \(L\) balls of volume
\(\sim\tau^D\) each, whereas the support near a landmark is a sheet of
extent \(\tau\) in \(m\) tangent directions and thickness \(t\ll\tau\) in the
remaining \(D-m\), i.e.\ volume \(\sim\tau^m t^{D-m}\). The ratio is
\((\tau/t)^{D-m}\): exponentially generous.

The fix is to split each neighbourhood by a local PCA chart. Let \(V_\ell\in
\R^{D\times m}\) be the top-\(m\) right singular vectors of the centred
training scatter in bucket \(\ell\), and for \(u=x-M_\ell\) write
\(u_\parallel=V_\ell^{\top}u\), \(u_\perp=u-V_\ell u_\parallel\). With
\(\sigma^\parallel_\ell,\sigma^\perp_\ell\) the median training magnitudes of
each part,
\begin{equation}
\label{eq:chartscore}
s(x)=\min_\ell
\sqrt{
\frac{\norm{u_\parallel}^2}{(\sigma^\parallel_\ell)^2}
+
\frac{\norm{u_\perp}^2}{(\sigma^\perp_\ell)^2}
}.
\end{equation}
This replaces balls with tangent-aligned ellipsoids. It is the \emph{default}
score; \eqref{eq:covlocal} is the fallback, and is recovered exactly by
setting \(V_\ell=0\) and
\(\sigma^\parallel_\ell=\sigma^\perp_\ell=r_\ell\), which is what buckets with
too few training points to estimate a chart fall back to automatically.

Measured on a 2-D sheet embedded in \(\R^8\), the ratio between the score of an
off-sheet displacement and an along-sheet displacement of equal length is:

\begin{center}
\begin{tabular}{lcc}
\toprule
sheet thickness \(t\) & locally-scaled balls & tangent charts \\
\midrule
\(0.05\)  & 1.34 & 2.3 \\
\(0.02\)  & 1.35 & 4.3 \\
\(0.005\) & 1.34 & 14.9 \\
\(0.001\) & 1.34 & 74.2 \\
\bottomrule
\end{tabular}
\end{center}

Balls are indifferent to thickness --- the discrimination sits at 1.34
regardless --- while the chart score scales as \(1/t\), which is the
\((\tau/t)^{D-m}\) argument showing up as a measurable quantity rather than an
asymptotic remark.

\subsection{\texorpdfstring{Conformal \(p\)-values}{Conformal p-values}}
\label{sec:pvalue}

With calibration scores \(\{s_j\}_{j=1}^{n_{\mathrm{cal}}}\) sorted
descending,
\begin{equation}
\label{eq:pvalue}
p(x)
=
\frac{1+\#\{s_j\ge s(x)\}}{n_{\mathrm{cal}}+1}.
\end{equation}
Small \(p\) means large cover (OOD). The resolution floor is
\(1/(n_{\mathrm{cal}}+1)\). Under exchangeability, \(p\) is super-uniform, so
flagging \(p\le\alpha\) controls the on-manifold false-alarm rate at level
\(\alpha\)~\citep{vovk2005}.

Batch tests use a permutation Mann--Whitney \(U\) on score samples; multiple
testing uses Benjamini--Hochberg.

\subsection{Closed-form minimum-norm repair}
\label{sec:repair}

The acceptance region \(\{\mathrm{cover}(x)\le\tau\}\) is a union of balls
\(\bigcup_\ell B(M_\ell,\tau r_\ell)\). Projection onto a union of balls is the
projection onto \emph{one} of them, but --- and this is where the earlier
statement of this section was wrong --- not in general the one with the nearest
centre. The travel required to reach ball \(\ell\) is \(d_\ell-\tau r_\ell\),
so the correct target is
\begin{equation}
\label{eq:repairstar}
\ell^\star=\argmin_\ell\ \bigl(d_\ell-\tau r_\ell\bigr),
\qquad
d_\ell=\norm{x-M_\ell},
\end{equation}
which coincides with \(\argmin_\ell d_\ell\) only when all \(r_\ell\) are
equal. A farther landmark with a more generous radius can be the cheaper
repair, and with the per-landmark radii of~\eqref{eq:covlocal} the radii are
never equal by design. Given \(\ell^\star\) the move is still closed form,
\begin{equation}
\label{eq:repair}
\delta^\star
=
\Bigl(1-\frac{\tau r_{\ell^\star}}{d_{\ell^\star}}\Bigr)(M_{\ell^\star}-x),
\end{equation}
equivalently the convex combination \(x_{\mathrm{repaired}}=
(\tau r_{\ell^\star}/d_{\ell^\star})\,x+
(1-\tau r_{\ell^\star}/d_{\ell^\star})\,M_{\ell^\star}\), and is a no-op when
\(x\) is already inside. The repair cannot invent structure outside the
landmark support. Because landmarks are a coarse support model, the certificate
can be loose relative to finer manifold geometry (e.g.\ EMD to real digits).

Under the chart score~\eqref{eq:chartscore} the target is selected the same
way and the point is moved onto that landmark's isotropic ball of radius
\(\tau r_{\ell^\star}\), a conservative inner move; the exact ellipsoid
projection reduces to a 2-D problem in
\((\norm{\delta_\parallel},\norm{\delta_\perp})\) and is not needed to certify
membership.

\subsection{Amortised negative space (CQR)}
\label{sec:negspace}

An optional quantile head predicts distance-to-support
\(y(\tilde x)=\min_j\norm{\tilde x-x_j}\) from FiLM features, trained with
pinball loss on isotropic perturbations and calibrated by conformalised
quantile regression~\citep{romano2019cqr}. Novelty detection calibrates only
the on-manifold null (raw median head output), remaining valid under covariate
shift of the negatives.

\begin{remark}[Why map-AUROC of unstructured probes misleads]
UMAP's \texttt{transform} places new points as convex combinations of training
neighbours' embeddings, so it cannot emit coordinates outside the convex hull
of the training map. Detectability of unstructured probes is therefore largely
about reserved empty space in the layout, not about representing ``far away.''
\end{remark}

%==============================================================================
\section{Measurements}
\label{sec:measurements}
%==============================================================================

Every number in this section is a mean over three seeds with the sample sd
reported, at \(120\) epochs on the shipping recipe, with a \(20\%\) holdout. A
difference is called \emph{resolved} only when it exceeds twice the pooled seed
sd of the two arms; otherwise it is inside run-to-run noise and is reported as
such. The seed spread is not a formality: on digits, \(\mathrm{5NN}\) accuracy
alone carries \(\mathrm{sd}=0.031\) across seeds, which is larger than most of
the differences this note previously drew conclusions from.

\subsection{Net effect of the changes}
\label{sec:regression}

A seeded snapshot was taken \emph{before} any of the changes described in this
note, so that a post-change regression could be told apart from a fix. Against
that baseline, on the shipping recipe, four differences are resolved out of
thirty metric--dataset pairs:

\begin{center}
\small
\begin{tabular}{llrr}
\toprule
Dataset & Metric & before & after \\
\midrule
digits     & density Spearman & \(0.717\pm0.020\) & \(0.670\pm0.007\) \\
swiss roll & \(k\)NN overlap  & \(0.793\pm0.005\) & \(0.840\pm0.004\) \\
swiss roll & area sd          & \(0.201\pm0.015\) & \(0.149\pm0.009\) \\
swiss roll & trustworthiness  & \(0.9953\pm0.0002\) & \(0.9972\pm0.0004\) \\
\bottomrule
\end{tabular}
\end{center}

Three are improvements on the manifold dataset and one is the known price of
the monotone squash (\S\ref{sec:squash}), at the magnitude that section
predicts. Everything else --- including all of the s-curve column, where
\(k\)-NN overlap rises by \(0.061\) and density Spearman by \(0.105\) but the
seed spread is too wide to resolve either --- is unchanged within noise. In
particular the FiLM reordering, the \(\varepsilon\) quantile estimator, and the
halo pass do not show up as regressions anywhere.

\subsection{Is FiLM conditioning earning its keep?}
\label{sec:filmvsconcat}

The landmark affinity \(a(x)\) is a deterministic function of \(x\), so
conditioning adds no information the encoder could not compute for itself. FiLM
is therefore an \emph{inductive bias} --- a soft partition-of-unity mixture of
experts --- and not extra capacity. The honest test is a plain MLP of identical
width, depth, and skip structure with \(a(x)\) simply \emph{concatenated} to
the normalised input.

\begin{center}
\small
\begin{tabular}{llcc}
\toprule
Dataset & Metric & FiLM & concat \\
\midrule
digits & 5-NN accuracy      & \(0.914\pm0.031\) & \(0.908\pm0.022\) \\
digits & trustworthiness    & \(0.937\pm0.010\) & \(0.934\pm0.006\) \\
digits & kNN overlap        & \(0.557\pm0.034\) & \(0.540\pm0.023\) \\
digits & geodesic Spearman  & \(0.724\pm0.060\) & \(0.629\pm0.024\)\;\textbf{*} \\
swiss roll & 5-NN accuracy  & \(0.922\pm0.024\) & \(0.918\pm0.029\) \\
swiss roll & kNN overlap    & \(0.860\pm0.005\) & \(0.850\pm0.009\) \\
swiss roll & geodesic Spearman & \(0.751\pm0.030\) & \(0.751\pm0.030\) \\
\bottomrule
\end{tabular}
\end{center}

\textbf{*} is the only resolved difference in the table. FiLM buys global
geodesic fidelity on clustered data --- \(+0.096\), roughly four times the
pooled spread --- and nothing else that is measurable. On the manifold dataset
the two are indistinguishable on every metric, to three decimals on geodesic
Spearman.

This is a narrower claim than the architecture's prominence suggests, and it is
the claim the evidence supports: the roles, temperatures, clamps, and
perplexity calibration are justified by one axis on one kind of data. Where
global structure is not the goal, concatenation is a legitimate and simpler
choice, and is available as \(\texttt{conditioning}=\texttt{"concat"}\).

\subsection{Persistence}
\label{sec:persistenceresults}

Following the protocol of \S\ref{sec:persistence}: three refits, Procrustes
fit on a shared anchor set and scored on held-out points.

\begin{center}
\small
\begin{tabular}{llccc}
\toprule
Dataset & Refit varies & rank Spearman & Jaccard@15 & coord.\ disagreement \\
\midrule
digits     & seed      & 0.712 & 0.377 & 0.729 \\
digits     & subsample & 0.633 & 0.377 & 0.726 \\
s-curve    & seed      & 0.998 & 0.727 & 0.032 \\
s-curve    & subsample & 0.995 & 0.657 & 0.264 \\
swiss roll & seed      & 0.998 & 0.716 & 0.030 \\
swiss roll & subsample & 0.862 & 0.609 & 0.724 \\
\bottomrule
\end{tabular}
\end{center}

Three things follow.

\emph{Manifold maps are reproducible; the clustered map is not.} At a fixed
training set, s-curve and swiss roll return the same map to rank Spearman
\(0.998\) and \(3\%\) coordinate disagreement. Digits reaches only \(0.712\),
with a Jaccard overlap of \(0.377\) --- under two neighbours in five are shared
between two runs that differ only in seed. Any downstream use of digits
coordinates as if they were stable is unsupported.

\emph{The s-curve subsample row is the gauge artefact of
Remark~\ref{rem:persprimary}, exactly as predicted.} Coordinate disagreement
rises eightfold (\(0.032\to0.264\)) while rank agreement is essentially
unchanged (\(0.998\to0.995\)). Reading the coordinate number alone would call
this a large instability; it is a reparameterisation.

\emph{Subsampling is the harder perturbation, and not uniformly so.} Swiss roll
degrades genuinely under subsampling (rank \(0.998\to0.862\)) rather than
merely moving, which is a real sensitivity of the unrolling to which points are
present.

\subsection{Ablations}
\label{sec:ablations}

\paragraph{Multiplicity \(\beta\).}
On digits, \(\beta\in\{0,0.5,1\}\) is inert: every metric differs by less than
\(0.01\) and none is resolved. This is the expected result and a useful
negative --- digits has no duplicate mass for the reweighting to act on, which
confirms \(\beta\) is a property of the data rather than a solver knob
(Remark~\ref{rem:beta}). It should be exercised where redundancy is real.

\paragraph{Pyramid against the geodesic anchor.}
The question was whether coarse-heavy level weights and
\(\lambda_{\mathrm{geo}}\) are substitutes. On swiss roll they are not, and the
result is stronger than a null:

\begin{center}
\small
\begin{tabular}{lccc}
\toprule
& \((1,1,1)\), \(\lambda_{\mathrm{geo}}{=}0.5\)
& \((1,2,8)\), \(\lambda_{\mathrm{geo}}{=}0.5\)
& \((1,2,8)\), \(\lambda_{\mathrm{geo}}{=}0\) \\
\midrule
\(k\)NN overlap & \(0.920\pm0.006\) & \(0.868\pm0.004\) & \(0.817\pm0.029\) \\
area sd         & \(0.081\pm0.004\) & \(0.124\pm0.006\) & \(0.181\pm0.040\) \\
\bottomrule
\end{tabular}
\end{center}

Removing the anchor at fixed \((1,2,8)\) costs \(0.051\) of \(k\)-NN overlap
and \(0.057\) of \(\mathrm{area\_sd}\), both resolved: the pyramid supplies
long-range attraction but not the global gauge, and cannot stand in for it.

The unexpected column is the first. At fixed \(\lambda_{\mathrm{geo}}=0.5\),
\emph{flat} level weights beat the coarse-heavy ramp by \(0.052\) of \(k\)-NN
overlap and \(0.043\) of \(\mathrm{area\_sd}\), both resolved. On a smooth
manifold with the geodesic anchor active, the ramp is not merely unnecessary
--- it is harmful, presumably because the anchor already carries global
structure and the coarse levels then compete with the fine ones for the edge
budget. The \((1,2,8)\) default is retained for the general case, where
\(\lambda_{\mathrm{geo}}=0.15\), but \S\ref{sec:defaults} now pairs raising
\(\lambda_{\mathrm{geo}}\) on smooth manifolds with flattening the level
weights, which was previously left as two independent knobs.

\paragraph{Geodesic anchor weight.}
On swiss roll with the frame term on, \(\lambda_{\mathrm{anchor}}=0\) is not
resolved from \(\lambda_{\mathrm{anchor}}=1\) on any metric, which supports the
view that \(\loss_{\mathrm{frame}}\) already carries local metric fidelity. It
is, however, markedly less \emph{stable}: \(\mathrm{area\_sd}\) is
\(0.131\pm0.072\) with the anchor off against \(0.075\pm0.008\) with it on, an
order of magnitude more seed spread. The default stays at \(1\); turning it off
is reasonable when the MDS diagnostic of \S\ref{sec:geo} is large, and should
be checked across seeds rather than one.

\paragraph{PCA skip and differential learning rates.}
On digits the differential rate does what \S\ref{sec:film} predicts, without
overturning the shipping recipe. Against the \(\texttt{pca\_skip}=
\texttt{False}\) control, the skip path at a flat rate loses \(0.098\) of
geodesic Spearman (resolved); at \(10\)--\(20\times\) on the head and
hypernetworks that loss falls to \(0.013\)--\(0.009\) and is no longer
resolved, and 5-NN accuracy recovers from \(0.851\) to \(0.871\).
Trustworthiness and kNN overlap remain resolved-worse than the no-skip control
in every arm, so the multiplier is a genuine improvement \emph{to the skip
path} rather than a reason to prefer it. Defaults are unchanged.

\subsection{Inference cost}
\label{sec:inference}

Out-of-sample speed is an explicit design goal, so it should carry numbers.
Digits (\(N=1797\), \(D=64\)), CPU, median of 15 repeats after warm-up, against
\texttt{UMAP.transform} and a PCA projection:

\begin{center}
\small
\begin{tabular}{lrrrr}
\toprule
& model on disk & \(B=1\) & \(B=512\) & \(\mu\)s / point at \(B=512\) \\
\midrule
leanmap & 944 KB & 0.25 ms & 1.73 ms & 3.4 \\
UMAP    & 605 KB & 3.22 ms & 399.9 ms & 781.0 \\
PCA     & 1 KB   & 0.03 ms & 0.05 ms & 0.09 \\
\bottomrule
\end{tabular}
\end{center}

leanmap is \(231\times\) faster than \texttt{UMAP.transform} per point at
\(B=512\) and \(13\times\) at \(B=1\); PCA is another \(38\times\) faster than
leanmap and is the honest floor --- a single \(D\times d\) matmul with no
conditioning --- included so the comparison is not read as leanmap being fast
in absolute terms.

\begin{remark}[What is actually eliminated]
\label{rem:notindexfree}
The neighbour search is not removed. It is \emph{replaced}: an exact
variable-size search against the training set becomes a fixed \(L\times D\)
dense distance computation against the landmarks, followed by a forward pass.
That is a different computation with a different answer, not a faster route to
the same one; the speed comes from having decided in advance how many
comparisons inference is allowed to make.

The trade is still the right one for the intended use, and for reasons that are
structural rather than constant-factor: no index to build or keep warm, no
training data shipped with the model, and a differentiable map that composes
into a downstream network. The last of these has no analogue in the UMAP
column at any speed.
\end{remark}

The artefact size difference is asymptotic, not the \(944\) versus \(605\) KB
the table shows at this \(N\). \texttt{UMAP.transform} needs the training data
and its neighbour graph retained, so the pickle grows linearly --- measured at
\(431\) bytes per point, flat to within \(1\%\) over \(N=250\) to \(1797\)
(\(106\to756\) KB). The leanmap artefact is \(L\times D\) landmarks plus
weights and is independent of \(N\). The crossover is near \(N\approx 2200\);
by \(N=10^6\) the UMAP artefact is roughly \(400\) MB against an unchanged
\(944\) KB.

%==============================================================================
\section{Default hyperparameters and interaction rules}
\label{sec:defaults}
%==============================================================================

Selected defaults for \(N\le 5\cdot 10^3\) via \texttt{PLANEConfig.for\_scale}:

\begin{center}
\small
\begin{tabular}{llc}
\toprule
Knob & Default (\(N\le 5\mathrm{k}\)) & Role \\
\midrule
\texttt{pca\_skip} & \texttt{False} & free residual layout \\
\texttt{lr} & \(2\cdot 10^{-2}\) & paired with skip off \\
\texttt{pca\_lr\_mult} & \(1\) & raise to \(10\)--\(20\) \emph{with} the skip \\
\texttt{min\_dist} & \(0.5\) & uniformity; see \S\ref{sec:bstability} \\
\(\lambda_{\mathrm{geo}}\) & \(0.15\) & \(0.5\) on smooth manifolds, with
flat weights \\
\(\lambda_{\mathrm{anchor}}\) & \(1\) & seed stability; \S\ref{sec:ablations} \\
\texttt{epochs} & \(240\) & --- \\
\texttt{pyramid\_level\_weights} & \((1,2,8)\) & coarse-heavy \\
\texttt{pyramid\_squash} & \texttt{rational\_q99} & monotone; \S\ref{sec:squash} \\
\texttt{conditioning} & \texttt{film} & see \S\ref{sec:filmvsconcat} \\
\texttt{width}/\texttt{depth} & \(384/3\) & capacity \\
\(\lambda_{\mathrm{frame}}\) & \(0\) & enable on fold-back sheets \\
\texttt{geo\_ramp} & \((0.2,0.45)\) & delayed gauge lock \\
\(\beta\) (\texttt{beta\_multiplicity}) & \(0.5\) & duplicate credit; see
Remark~\ref{rem:beta} \\
\bottomrule
\end{tabular}
\end{center}

Fixed plumbing (not exposed as knobs):
\(\lambda_{\mathrm{backbone}}=0.01\), \(\lambda_{\mathrm{ord}}=0.1\),
\(\eta=1\), \(\mathrm{spread}=1\), \(\texttt{FRAME\_CENTERS}=128\),
\(\texttt{FRAME\_NORMAL\_THRESH}=0.5\), \(\gamma\in[0.1,10]\).

Three non-separable design rules:

\begin{enumerate}[leftmargin=*]
\item \texttt{pca\_skip}, \(\mathrm{lr}\), and \texttt{pca\_lr\_mult} are one
  decision (Section~\ref{sec:film}). The escape from the shared-rate conflict
  is the parameter group, not a compromise rate.
\item Raising \texttt{min\_dist} buys uniformity at no measured cost in label
  accuracy; \(0.5\) captures nearly all of the available gain
  (Remark~\ref{rem:mindistseeded}). The former ``never below \(0.2\)'' rule is
  withdrawn.
\item Fold-back manifolds need delayed frame rigidity
  (Section~\ref{sec:frame}).
\item \(\lambda_{\mathrm{geo}}\) and \texttt{pyramid\_level\_weights} are one
  decision, not two. With the anchor strong on a smooth manifold, flat weights
  are resolved-better than the coarse-heavy ramp (\S\ref{sec:ablations}).
\end{enumerate}

Defaults \emph{not} changed by this pass, and why: \(\beta\) is inert on the
data measured and is a property of the corpus rather than the solver;
\(\lambda_{\mathrm{anchor}}\) is unresolved on the mean but halves the seed
spread; \texttt{conditioning} stays on FiLM for the one axis where it is
resolved. In each case the evidence is reported in \S\ref{sec:measurements}
rather than compiled into a silent adaptive rule.

%==============================================================================
\section{Relation to prior methods}
\label{sec:related}
%==============================================================================

leanmap sits at the intersection of several classical threads:

\begin{itemize}[leftmargin=*]
\item \textbf{Fuzzy neighbour embedding.} The geometric loss and \((a,b)\)
  kernel follow UMAP~\citep{mcinnes2018umap}; the graph uses the same
  \(\sigma/\rho\) solve and fuzzy union. Differences: parametric \(f_\theta\),
  direct autodiff (hence the raised \(\texttt{min\_dist}\)), multi-scale
  Galerkin pyramid, and metric regularisers.
\item \textbf{Isomap / classical MDS.} Landmark geodesics and
  \(\loss_{\mathrm{anchor}}\) import the Isomap gauge~\citep{isomap2000} as a
  soft prior rather than as the embedding itself.
\item \textbf{Local tangent space alignment.} In tangent mode
  \(\loss_{\mathrm{frame}}\) is closest to LTSA~\citep{zhang2004ltsa}: both
  estimate a weighted local tangent basis at each neighbourhood and require the
  structure expressed in those tangent coordinates to agree between the ambient
  data and the embedding. The differences are that leanmap matches Gram
  matrices rather than aligning tangent coordinates through a global
  eigenproblem, detaches the similarity scale so the objective is scale-free,
  and applies the constraint as a penalty on a parametric map instead of
  solving for the embedding in closed form. LTSA is the closer ancestor; ARAP
  is the secondary reading.
\item \textbf{As-rigid-as-possible deformation.} \(\loss_{\mathrm{frame}}\) can
  equally be read as ARAP~\citep{sorkine2007arap} expressed through neighbour
  Gram matrices with a detached similarity scale.
\item \textbf{Parametric UMAP / NN embedding.} The FiLM-conditioned MLP is a
  parametric embedder in the spirit of parametric
  UMAP~\citep{sainburg2021parametric}, with landmark affinities as the
  conditioning code and an optional PCA residual path.
\item \textbf{Conformal prediction.} Cover scores yield finite-sample,
  distribution-free OOD control under exchangeability~\citep{vovk2005}.
\end{itemize}

%==============================================================================
\section*{Appendix: loss assembly (implementation mirror)}
%==============================================================================

\begin{sloppypar}
The reference trainer assembles~\eqref{eq:total} each step from the modules
\texttt{fuzzy\_cross\_entropy}, \texttt{ordinal\_triplet\_loss},
\texttt{landmark\_regularisation}, \texttt{local\_rigidity\_loss},
\texttt{procrustes\_anchor\_loss}, and \texttt{geodesic\_stress\_loss} in
\path{src/leanmap/losses.py}, with ramps from \texttt{alignment\_ramp}. Graph
stages live in \path{src/leanmap/graph.py}; landmarks and classical MDS in
\path{src/leanmap/landmarks.py}; the encoder in \path{src/leanmap/model.py};
conformal scores in \path{src/leanmap/conformal.py}.
\end{sloppypar}

\begin{thebibliography}{99}

\bibitem{mcinnes2018umap}
L.~McInnes, J.~Healy, and J.~Melville.
\newblock UMAP: Uniform Manifold Approximation and Projection for Dimension
Reduction.
\newblock \emph{arXiv:1802.03426}, 2018.

\bibitem{isomap2000}
J.~B.~Tenenbaum, V.~de~Silva, and J.~C.~Langford.
\newblock A global geometric framework for nonlinear dimensionality reduction.
\newblock \emph{Science}, 290(5500):2319--2323, 2000.

\bibitem{torgerson1952}
W.~S.~Torgerson.
\newblock Multidimensional scaling: I.\ Theory and method.
\newblock \emph{Psychometrika}, 17(4):401--419, 1952.

\bibitem{sorkine2007arap}
O.~Sorkine and M.~Alexa.
\newblock As-Rigid-As-Possible Surface Modeling.
\newblock In \emph{Symposium on Geometry Processing}, 2007.

\bibitem{zhang2004ltsa}
Z.~Zhang and H.~Zha.
\newblock Principal manifolds and nonlinear dimensionality reduction via
tangent space alignment.
\newblock \emph{SIAM Journal on Scientific Computing}, 26(1):313--338, 2004.

\bibitem{sainburg2021parametric}
T.~Sainburg, L.~McInnes, and T.~Q.~Gentner.
\newblock Parametric UMAP embeddings for representation and semi-supervised
learning.
\newblock \emph{Neural Computation}, 33(11):2881--2907, 2021.

\bibitem{levina2005}
E.~Levina and P.~J.~Bickel.
\newblock Maximum likelihood estimation of intrinsic dimension.
\newblock In \emph{NeurIPS}, 2005.

\bibitem{vovk2005}
V.~Vovk, A.~Gammerman, and G.~Shafer.
\newblock \emph{Algorithmic Learning in a Random World}.
\newblock Springer, 2005.

\bibitem{romano2019cqr}
Y.~Romano, E.~Patterson, and E.~J.~Cand\`es.
\newblock Conformalized Quantile Regression.
\newblock In \emph{NeurIPS}, 2019.

\bibitem{leanmap-config}
leanmap documentation.
\newblock Configuration guide.
\newblock \texttt{docs/CONFIGURATION.md}, 2026.

\end{thebibliography}

\end{document}
